\scanpage{4}
\mspage{2}
\deleted{\unsure{s'immerge injectivement} dans un groupe algébrique
commutatif de rang unipotent nul
(i.e. ext. d'une VA par un tore) alors
\uline{l'application}}
\deleted{\diagram{\begin{tikzcd}[ampersand replacement=\&]
X(\overline{\mathbf Q})\arrow[r]\&{}
\end{tikzcd}}}

\textcircled{3} Supposons $X$ de t.f. sur $\overline{\mathbf Q}$. Alors on a
un ``germe de relèvement'' de $\Gamma$
\[
(9)\qquad\diagram{\begin{tikzcd}[ampersand replacement=\&]
\mathcal G_{\overline X}\arrow[r]\&\Gamma\arrow[d,phantom,"\cup" description]\\
\&\Gamma'\arrow[ul]
\end{tikzcd}}
\]
sur ss-groupe ouvert (modulo la relation
de coïncidence sur un ss groupe
ouvert plus petit). Pour \inserted{$\xi\in X(\overline{\mathbf Q})$ donné, et pour}
$\Gamma'$ assez petit,
$\Gamma'\mathbin{-}\mathcal G_{\overline X}$ se factorise par $\mathcal G_{\overline X}(\xi)$, de
\margin{$X=\overline X$}
sorte qu'on a un \inserted{germe d'} hom composé
\[
\deleted{\Gamma'\longrightarrow\mathcal E_X}
\qquad\inserted{\Gamma'\longrightarrow\mathcal E_X}
\]
rendant commutatif
\[
(10)\qquad\diagram{\begin{tikzcd}[ampersand replacement=\&]
\mathcal E_X\arrow[r]\&\mathcal G_X\arrow[r]\&\Gamma\\
\&\&\Gamma'\arrow[u,hook]\arrow[ull,"{\unsure{\varphi_\xi^t|\varphi_\xi}}"']
\end{tikzcd}}
\]
-- du moins si $\widetilde X=\widetilde X(\xi)$ -- sinon un
hom défini mod composition par un \unsure{autom}
int. provenant de $\pi_1(X)$. \unsure{Soulignons}
\mspage{3}
